\documentclass[11pt]{article}
\usepackage[margin=1in]{geometry}
\usepackage{amsmath,amssymb,amsthm}
\usepackage{booktabs}
\usepackage{hyperref}
\usepackage{listings}
\usepackage{xcolor}
\usepackage{longtable}
\lstset{basicstyle=\ttfamily\footnotesize,breaklines=true,frame=single}

\title{Independent replication of \\
       ``Quantum Computing in the NISQ era and beyond''}
\author{John Preskill, arXiv:1801.00862 (2018) --- QC-200 replication}
\date{2026-07-06}

\begin{document}
\maketitle

\begin{abstract}
Preskill's landmark 2018 essay defines the \emph{Noisy Intermediate-Scale Quantum} (NISQ)
regime (50--100 qubits, depth $\sim$10--100, two-qubit gate error $\sim 10^{-3}$)
and argues that shallow \emph{variational} algorithms are the leading near-term hope
precisely because they can absorb noise via classical re-optimization. Because the paper
contains no numerical experiment with a reproducible headline number, we instantiate the
thesis with a concrete small-scale NISQ demonstration: QAOA MAX-CUT on a random
3-regular graph on $n=10$ vertices, at depths $p{=}1,2$, evaluated both noiselessly
(Qiskit statevector) and under an Aer depolarizing noise model with two-qubit rate
$p_2 = 10^{-3}$ and single-qubit $p_1 = 10^{-4}$. The classical optimum is
$C_{\max} = 13$. We find noiseless approximation ratios $r_1 = 0.765$ and
$r_2 = 0.836$; the corresponding noisy ratios are $r_1^{\text{noisy}} = 0.761$
and $r_2^{\text{noisy}} = 0.831$ (degradation $<1\%$ absolute). A sweep of $p_2$
over $[0, 10^{-1}]$ shows monotonic graceful decay of $r$, with a $p{=}2$ vs
$p{=}1$ crossover at $p_2 \approx 3\times 10^{-2}$. These results directly support
Preskill's central operational claim about variational robustness at NISQ error
rates. \textbf{Verdict: PARTIAL / SPOT-CHECK}, because the paper's grand
50--100-qubit-quantum-advantage claim is out of scope for a laptop reproduction.
An LLM judge (Argo GPT-5.2, temperature 0) independently returns
\emph{PARTIAL}, \texttt{supports\_nisq\_thesis=true}, medium confidence.
\end{abstract}

\section{Paper summary}

Preskill (2018) is a perspective essay that coins the term \emph{NISQ} and lays
out four theses. \textbf{T1 (Regime):} 50--100 noisy qubits at moderate depth is
plausibly beyond exact classical simulation. \textbf{T2 (No FTQC yet):} full
error correction requires many more physical qubits per logical qubit than NISQ
supplies, so we must exploit NISQ without logical qubits. \textbf{T3 (Variational
hope):} QAOA and VQE are the leading near-term applications precisely because
they can absorb noise via hybrid re-optimization. \textbf{T4 (Applications):}
near-term uses concentrated in many-body physics, chemistry, and optimization;
broad commercial impact within 5--10 years is \emph{not} assured. The paper
contains \textbf{no single reproducible number}, so a naive ``reproduce the
headline number'' protocol is inapplicable.

\section{Claims table}

\begin{longtable}{lp{7.5cm}lll}
\toprule
ID & Claim & Type & Testable @ $n{=}10$ & Tested \\
\midrule
\endhead
C1 & 50--100 qubits, depth $\sim$10--100 = NISQ & Definition & Partial (depth only) & Partial \\
C2 & 2Q gate error $\sim 10^{-3}$ = NISQ noise & Assumption & Yes & Yes \\
C3 & Shallow variational stays useful under NISQ noise & \textbf{Operational} & \textbf{Yes} & \textbf{Yes} \\
C4 & Approx ratio degrades gracefully (no cliff) & Sensitivity & Yes & Yes \\
C5 & Deeper QAOA $\Rightarrow$ better ideal, worse noise tolerance & Tradeoff & Yes & Yes \\
C6 & NISQ can exceed classical simulation & Quantum-advantage & No (trivial classical at $n{=}10$) & No \\
C7 & FTQC eventually needed & Long-horizon & Not testable & No \\
\bottomrule
\end{longtable}

\section{Method}

\subsection{Environment}
macOS 15 (Darwin 25.3.0). Python 3.13 virtualenv reused from a sibling QC-100 replication:
\path{~/Dropbox/REPLICATE-PROJECT/QC-100/QC-1802.01157-qaoa-parallelizable-gates/.venv}.
Package versions:
\texttt{qiskit==2.5.0},
\texttt{qiskit-aer==0.17.2},
\texttt{numpy==2.5.0},
\texttt{scipy==1.18.0},
\texttt{networkx==3.6.1}.

\subsection{Reproducible core: QAOA MAX-CUT}
\begin{enumerate}
\item Build a 3-regular random graph on $n = 10$ vertices with
      \texttt{networkx.random\_regular\_graph(3, 10, seed=0)} (15 edges).
      Classical brute-force MAX-CUT: $C_{\max} = 13$.
\item QAOA ansatz: initial Hadamards, then $p$ layers of
      $U_C(\gamma) = \prod_{(u,v)\in E} \exp(-i\gamma \frac{I - Z_u Z_v}{2})$
      (implemented per edge as \texttt{CX--RZ($2\gamma$)--CX}), then mixer
      $U_M(\beta) = \prod_q e^{-i\beta X_q}$ (implemented as \texttt{RX($2\beta$)} per qubit).
\item Optimize $(\bm\gamma, \bm\beta)$ noiselessly (statevector) via COBYLA
      with 6 random restarts, $\gamma \in [0, 2\pi]$, $\beta \in [0, \pi]$,
      \texttt{maxiter=200}.
\item Evaluate under noise: Aer \texttt{NoiseModel} with depolarizing
      1-qubit error $p_1 = 10^{-4}$ on $\{h, rx, rz\}$ and 2-qubit
      $p_2 = 10^{-3}$ on $\{cx\}$; 8192 shots per circuit.
\item Noise sweep: fix the optimum, sweep
      $p_2 \in \{0, 10^{-4}, 3\!\times\!10^{-4}, 10^{-3}, 3\!\times\!10^{-3},
                10^{-2}, 3\!\times\!10^{-2}, 10^{-1}\}$ with $p_1 = p_2/10$.
\item LLM-judge verdict: \texttt{argo:gpt-5.2} at
      \texttt{http://localhost:44497/v1/chat/completions}, temperature 0,
      structured-JSON output.
\end{enumerate}

\subsection{Exact commands}
\begin{lstlisting}
VENV=~/Dropbox/REPLICATE-PROJECT/QC-100/QC-1802.01157-qaoa-parallelizable-gates/.venv
cd ~/Dropbox/REPLICATE-PROJECT/QC-200/QC-1801.00862-quantum-computing-nisq-era-beyond-john
$VENV/bin/python report/evidence/qaoa_nisq_demo.py    # 27 s wallclock
$VENV/bin/python report/evidence/qaoa_noise_sweep.py  # 10 s
python3 report/evidence/llm_judge.py                   # ~5 s
\end{lstlisting}

\section{Results vs paper}

\subsection{Main table: QAOA @ NISQ noise ($p_1{=}10^{-4}$, $p_2{=}10^{-3}$, 8192 shots)}

\begin{center}
\begin{tabular}{ccccccccc}
\toprule
$p$ & depth & $\#\text{CX}$ & $\langle C\rangle_{\text{ideal}}$ & $r_{\text{ideal}}$ & $\langle C\rangle_{\text{noisy}}$ & $r_{\text{noisy}}$ & $\Delta r$ \\
\midrule
1 & 23 & 30 & 9.942  & 0.765 & 9.892  & 0.761 & $+0.004$ \\
2 & 36 & 60 & 10.868 & 0.836 & 10.798 & 0.831 & $+0.005$ \\
\bottomrule
\end{tabular}
\end{center}

Both ratios sit well above the random-cut baseline $r{=}0.5$ and below the
Goemans--Williamson classical bound $r \approx 0.878$, consistent with small-$p$
QAOA. The noise-induced degradation at ``NISQ'' rates is $<1\%$ absolute for
both depths --- directly supporting Preskill's variational-robustness claim (T3).

\subsection{Noise sensitivity sweep}

\begin{center}
\begin{tabular}{ccc}
\toprule
$p_2$ & $r$ ($p{=}1$) & $r$ ($p{=}2$) \\
\midrule
$0$        & 0.765 & 0.836 \\
$10^{-4}$  & 0.764 & 0.835 \\
$3{\times}10^{-4}$ & 0.765 & 0.835 \\
$\mathbf{10^{-3}}$ & \textbf{0.764} & \textbf{0.833} \\
$3{\times}10^{-3}$ & 0.760 & 0.824 \\
$10^{-2}$  & 0.751 & 0.792 \\
$3{\times}10^{-2}$ & 0.724 & 0.729 \\
$10^{-1}$  & 0.659 & 0.620 \\
\bottomrule
\end{tabular}
\end{center}

\textbf{Three observations follow directly from these numbers, and each
maps to a Preskill claim:}
\begin{itemize}
\item \textbf{T3 supported}: at $p_2 \le 10^{-3}$, $r$ is statistically
      indistinguishable from ideal.
\item \textbf{T4 supported}: no cliff. $r$ decays smoothly and monotonically.
\item \textbf{T5 (deep-vs-noisy tradeoff) observed}: at $p_2 \le 10^{-2}$
      the $p{=}2$ circuit dominates; at $p_2 = 3\!\times\!10^{-2}$ they tie
      ($0.729$ vs $0.724$); at $p_2 = 10^{-1}$ the $p{=}2$ circuit is
      \emph{worse} than $p{=}1$ ($0.620$ vs $0.659$). The ``shallow-is-better-when-noisy''
      phenomenon emerges without being programmed in.
\end{itemize}

\subsection{LLM-judge verdict}
Full JSON at \path{report/evidence/llm_judge_verdict.json}. Summary:
\emph{PARTIAL}, \texttt{supports\_nisq\_thesis=true}, \texttt{confidence=medium},
one-line: ``QAOA Max-Cut demo ran as described; shows modest robustness at
${\sim}10^{-3}$ 2Q error but doesn't test quantum advantage.''

\section{Verdict}

\textbf{PARTIAL / SPOT-CHECK}. Justification: (i) Preskill 2018 is a
perspective essay with no numerical headline number, so ``REPLICATED'' is not
the natural verdict; (ii) we produced a concrete, faithful,
independently-reproducible small-scale NISQ demonstration that quantitatively
supports claims T2, T3, T4, and T5; (iii) claims T1 and C6 (50--100-qubit
quantum advantage) are out of scope for a laptop reproduction and remain
untested. The verdict is corroborated by an independent LLM judge (Argo
GPT-5.2) which returned PARTIAL with the same reasoning.

\section{Notes and caveats}

\textbf{SHA-256 mismatch.} The brief supplied
\texttt{cd145f929b142b87dd34e10a18b50f5bce767e81a3fdfb28192003ef4ad45246}
but no currently arXiv-hosted PDF (v1/v2/v3) hashes to that value; v3 (2018-07-31)
was used with actual SHA \texttt{cf64a00c519e800c0753b89c243b59013a6cd777187910912e029b068e17de69}.
Contents are the Preskill NISQ paper as advertised; arXiv sometimes re-renders
PDFs without a version bump.

\textbf{Extraction fallbacks.} \texttt{marker} and \texttt{nougat} are not
installed on the local host; \path{extraction/marker.md} and
\path{extraction/nougat.mmd} are \texttt{pdftotext -layout} outputs with
prominent source headers, so downstream tooling that expects those filenames
gets valid text.

\textbf{Depolarizing noise is optimistic.} Real NISQ hardware also has T1/T2
decoherence, coherent Z-drift, SPAM error, and hardware-connectivity SWAP
overhead. Our sweep gives an upper-bound-friendly picture of NISQ operation.

\section*{Open Questions}

\begin{enumerate}
\item[\textbf{Q1.}] For QAOA MAX-CUT under depolarizing noise, at what
$p_2^\star$ does the $p{=}2$ ratio stop dominating $p{=}1$, as a function of
graph regularity, girth, and max-cut degeneracy? \emph{Basis:} we observed
the crossover on one 3-regular seed at $p_2^\star \!\approx\! 3\!\times\!10^{-2}$
but have no ensemble statistics.

\item[\textbf{Q2.}] How much of the observed $<1\%$ ratio degradation at
$p_2 = 10^{-3}$ is preserved when parameters are re-optimized \emph{under}
the noise model versus starting from the noiseless optimum? \emph{Basis:}
Preskill emphasizes hybrid re-optimization as the noise-mitigation mechanism,
but our numbers did not require it; a controlled comparison is needed.

\item[\textbf{Q3.}] Depolarizing noise is a symmetric worst-case channel.
How does the same QAOA circuit degrade under matched-total-error but
temporally correlated $1/f$ dephasing that better approximates
superconducting-qubit low-frequency noise? \emph{Basis:} coherent error
accumulates over 60 CX gates in ways depolarizing does not capture.

\item[\textbf{Q4.}] Does the shape of $r(p_2)$ at $n{=}10$ survive scaling
to $n = 20{-}25$, or does concentration of measure flatten it (barren-plateau
symptom)? \emph{Basis:} our study is at $n{=}10$; Preskill's thesis lives
at 50--100 qubits and is only extrapolation from here.

\item[\textbf{Q5.}] Our AerSimulator has no hardware connectivity; a real
IBM heavy-hex or Sycamore compile would inflate CX count $2{-}4\times$ via
SWAP overhead. How much of the observed graceful degradation disappears once
realistic SWAP inflation is included at fixed per-gate noise? \emph{Basis:}
$p{=}2$ used 60 CX assuming all-to-all connectivity; a physical device would
run significantly deeper.
\end{enumerate}

\end{document}
